% Generated draft: requires mathematical and source review.
\documentclass[11pt]{article}
\usepackage[margin=25mm]{geometry}
\usepackage{fontspec}
\setmainfont{Latin Modern Roman}
\usepackage{amsmath,amssymb,mathtools,unicode-math}
\setmathfont{Latin Modern Math}
\usepackage{enumitem,needspace,xurl,hyperref}
\setlength{\parindent}{0pt}
\setlength{\parskip}{5pt}
\setlength{\emergencystretch}{3em}
\newcommand{\N}{\mathbb{N}}
\newcommand{\Z}{\mathbb{Z}}
\newcommand{\Q}{\mathbb{Q}}
\newcommand{\R}{\mathbb{R}}
\newcommand{\C}{\mathbb{C}}
\newcommand{\F}{\mathbb{F}}
\newcommand{\A}{\mathbb{A}}
\newcommand{\PP}{\mathbb{P}}
\newcommand{\E}{\mathbb{E}}
\newcommand{\eps}{\varepsilon}
\newcommand{\dd}{\,\mathrm{d}}
\begin{document}
\section*{UnsolvedMath ID: 800723}
\textbf{Source problem code:} catalogue entry
\subsection{Definitions and mathematical statement}
A finite group \(H\) is \emph{metacyclic} if it has a cyclic normal subgroup \(C\trianglelefteq H\) such that \(H/C\) is cyclic. A finite group \(G\) is \emph{Sylow-metacyclic} if, for every prime \(p\mid |G|\), every Sylow \(p\)-subgroup of \(G\) is metacyclic.

For a field \(F\), a finite group \(G\) is \emph{\(F\)-admissible} if there exist a finite-dimensional central division algebra \(D\) over \(F\) and a maximal commutative subfield \(L\subseteq D\), containing \(F\), such that \(L/F\) is Galois and \(\operatorname{Gal}(L/F)\cong G\). Here central means \(Z(D)=F\), and maximal means that no larger commutative subfield of \(D\) contains \(L\). In this situation,
\[
[L:F]=\deg(D)=|G|,\qquad \dim_F D=|G|^2.
\]

\textbf{Statement (the \(\Q\)-admissibility conjecture).} For every finite Sylow-metacyclic group \(G\), there exist a finite-dimensional central division algebra \(D\) over \(\Q\) and a maximal subfield \(L\subseteq D\) such that \(L/\Q\) is Galois and
\[
\operatorname{Gal}(L/\Q)\cong G.
\]

The supplied extract describes the surrounding problem as an extension of the inverse Galois problem, especially interesting over \(\Q\), and mentions Problem~5 of the reference labeled \texttt{[Am:surv]}. Because this contextual passage precedes the words introducing the ``next problem,'' its precise connection to the displayed conjecture is unresolved. The extract says that the conjecture was first recognized in the reference labeled \texttt{[Schacher:th1]}; neither reference label is bibliographically resolved in the supplied record.
\subsection{Short English statement}
Every finite group whose Sylow subgroups each have a cyclic normal subgroup with cyclic quotient should occur as the Galois group of a maximal subfield of a finite-dimensional central division algebra over \(\Q\). The supplied registry labels this conjecture open; that status has not been independently verified.
\subsection{Sources}
\begin{itemize}
\item Local UnsolvedMath registry, permanent ID \(800723\), exact source problem code \texttt{catalogue entry}, title \emph{Open Problems on Central Simple Algebras --- Problem 38}. The record lists Auel, Brussel, Garibaldi, and Vishne under proposer attribution and gives the year 2010. Its supplied status is \texttt{open}.
\item Auel, Brussel, Garibaldi, and Vishne, \emph{Open Problems on Central Simple Algebras} (2010), Problem~38, \url{https://arxiv.org/abs/1006.3304}. The supplied provenance records access on 2026-10-16 and extraction by \texttt{arxiv\_latex}; these are supplied provenance details, not an independently verified visit.
\item Reference label \texttt{[Am:surv]}: the supplied contextual passage mentions Problem~5. No corresponding bibliographic entry is supplied.
\item Reference label \texttt{[Schacher:th1]}: the supplied extract attributes the first recognition of the \(\Q\)-admissibility conjecture to this reference. No corresponding bibliographic entry is supplied.
\end{itemize}
\end{document}
